% SPDX-License-Identifier: Apache-2.0
\section{Two Processes, One per Edge}
\label{sec:x-twoproc}

A unit may have several processes: \code{run} is \code{join2} of two
loops, each with a wait of its own, and each is a process in the
runtime and a clocked block in the netlist. The sampler here counts
on the rising edge while enabled and copies the count on the falling
edge, half a cycle later; the output shows the copy. The two loops
share the unit through a plain reference, \code{this}, since two
closures cannot both hold \code{\&mut self}, and the lowering reads
\code{this.field} as it reads \code{self.field}. The falling-edge
wait, \code{C::falling()}, lowers to a block on the negative edge in
Verilog and \code{falling\_edge} in VHDL; a register such a block
drives is listed as one in the ports sidecar, and the testbench checks
it after the falling edge rather than the rising one. The lowering is
simulated against the trace under both simulators,
Section~\ref{sec:x-checked}.

\begin{figure}[htbp]
\centering
\input{twoproc_timing.tex}
\caption{Two processes: \code{n} advances on the rising edge, \code{half}
copies it on the falling one, and \code{q} follows \code{half}.}
\label{fig:x-twoproc}
\end{figure}

\lstinputlisting[caption={\code{ex\_twoproc.rs}. Run with \code{bazel run //lib/examples:ex\_twoproc}.},label={lst:x-twoproc}]{lib/examples/ex_twoproc.rs}

\lstinputlisting[style=ascii,caption={What \code{ex\_twoproc} printed when the build ran it: the output per cycle, then the lowering with its two blocks.},label={lst:x-twoproc-out}]{out_twoproc.txt}
